\section{Introduction}\label{sec:intro}

A smishing message (from SMS phishing) imitates a bank, a delivery service or a government agency and leads the recipient to a link, a phone number or an application, with the aim of obtaining credentials or a payment. Unlike advertising spam, a missed smishing message can lead to direct financial loss, so the two kinds of error have unequal costs. Automatic detection works under three constraints. The messages are short, informal and often written in several languages. Senders alter spelling, insert spaces or replace letters with visually similar symbols in order to evade filters \cite{salman2024,boucher2022}. A filter that runs on a handset or at an operator gateway must classify every message within milliseconds on modest hardware.

Many studies report high accuracy on public corpora of short messages \cite{mishra2023dsmishsms,mishra2022detector,sonowal2020,roy2020,ghourabi2023,mahmud2024,shim2024}. Three aspects of the usual evaluation make such results difficult to compare and to transfer to practice.

\emph{Duplicates and overlap between corpora.} The Mendeley SMS Phishing Dataset \cite{mendeley2022,mishra2023dataset} is one of the few public corpora that separates smishing from spam. Of its \val{men-raw} messages, our deduplication removed \val{men-removed}: exact or near duplicates of a kept message and members of \val{men-conflict-groups} groups of identical or nearly identical texts that carry different labels. Of the \val{men-clean} remaining messages, \val{men-in-uci} (\val{men-in-uci-pct}\%) have an exact or near duplicate in the older UCI SMS Spam Collection \cite{almeida2011}. A random split of such data can place copies of one message in both the training and the test part, which is a known source of leakage and of optimistic estimates \cite{kapoor2023,lee2022}. The published methods that we re-implement do not state whether duplicates were removed \cite{mahmud2024,shim2024}.

\emph{Robustness to obfuscation.} Character-level perturbations are known to degrade word-level and subword neural models \cite{gao2018,li2019,pruthi2019}, and Unicode look-alike characters can change the input of a model without changing what a reader sees \cite{boucher2022}. When robustness is measured, the test perturbations are often of the same kind as those used in training. An evaluation should also include perturbation types that the model has never seen.

\emph{Inference cost.} Fine-tuned transformer encoders \cite{devlin2019,sanh2019,liu2019} are strong text classifiers, but in our measurements a DistilBERT classifier needs a median of \val{lat-distilaug-p50}~ms per message on one CPU thread, while a linear model over character n-grams needs \val{lat-svmcharaug-p50}~ms. Latency is rarely reported together with accuracy, and almost never end to end for a complete detection pipeline.

This paper studies whether a two-stage cascade can combine the speed and robustness of a character-level linear model with the accuracy of a transformer. The first stage is a linear support vector machine over TF-IDF weighted character n-grams \cite{cortes1995,joachims1998} with sigmoid calibration of its scores \cite{niculescu2005}. When the largest calibrated class probability of the first stage falls below a threshold $\tau$, the message is passed to the second stage, a DistilBERT classifier \cite{sanh2019}. Both stages are trained with random character obfuscation, and $\tau$ is chosen on validation data by an explicit criterion. The cascade is compared with classical, neural and transformer baselines, with a few-shot large language model \cite{qwen2024}, and with re-implementations of two recent published methods \cite{mahmud2024,shim2024} in their own task settings. All comparisons follow a protocol that was frozen before the test data were read.

The contributions are as follows.
\begin{enumerate}
\item An evaluation protocol for the detection of fraudulent messages: deduplicated splits of three corpora in English and Portuguese, checked for leakage between splits; a configuration frozen before testing; \val{n-seeds} training seeds; paired bootstrap and McNemar tests with Holm correction; and re-implementation of published methods in the task settings of their papers.
\item A formal description of the method: the cascade as a decision rule with a confidence threshold, the choice of the threshold as an optimisation problem on validation data, a cost function with unequal error costs, and obfuscation as a probabilistic string transformation with one family of operations used for training and another family held out from training.
\item Results reported together with their statistical status. On the clean three-class test set the cascade reaches a macro F1 of \val{cascade-clean-f1}~$\pm$~\val{cascade-clean-f1-sd} (mean $\pm$ standard deviation over seeds), and the character model alone without augmentation reaches \val{svmchar-clean-f1}; the difference is not significant (Holm-adjusted $p=\val{clean-best-p-boot-holm}$ for the bootstrap test and $p=\val{clean-best-p-mcn-holm}$ for the McNemar test). Under the strongest held-out obfuscation in the published task settings, the cascade keeps a macro F1 of \val{cascade-mahmud-ho20} and \val{cascade-shim-ho20}, while the re-implemented published methods fall to between \val{pub-ho20-min} and \val{pub-ho20-max}; these differences are significant. A post-hoc comparison shows that the character model alone is also significantly better than the published methods under this obfuscation. In the three-class task the cascade passes between \val{routed-min-pct}\% and \val{routed-max-pct}\% of the messages to the transformer and has a median end-to-end latency of \val{lat-cascade-p50}~ms on one CPU thread. In two of the binary settings the threshold chosen on validation data is high, and between \val{routed-binary-min-pct}\% and \val{routed-binary-max-pct}\% of the messages reach the transformer.
\item A qualitative case study for Kazakhstan with author-written messages in Russian, Kazakh and English, machine translation \cite{nllb2024}, adaptation of the character model with a small synthetic set, and four public real messages. The case study shows failure modes of models trained on English data; it is not a quantitative evaluation.
\end{enumerate}

The main practical finding is therefore modest. A calibrated character-level linear model is a strong, fast and robust detector, and in the three-class task the transformer stage of the cascade changes few of its decisions. Its role is a fallback for the small share of messages on which the first stage is not confident.

Section~\ref{sec:related} reviews related work. Section~\ref{sec:methods} describes the data, the obfuscation model, the cascade, the cost function, the statistical tests and the experimental protocol. Section~\ref{sec:results} presents the results, their analysis, the Kazakhstan case study and the limitations. Section~\ref{sec:conclusion} concludes.

\subsection{Related work}\label{sec:related}

\emph{Detection of SMS spam and smishing.} The UCI SMS Spam Collection \cite{almeida2011} has long been the standard benchmark for SMS spam filtering. Systems built specifically for smishing often combine content features with checks of the links and domains in a message. DSmishSMS \cite{mishra2023dsmishsms} and Smishing Detector \cite{mishra2022detector} use hand-crafted features together with external domain and URL checks, and Sonowal \cite{sonowal2020} ranks content features with several correlation measures before classification. Such pipelines depend on external services and on features that are not fully specified, which makes exact re-implementation difficult. Neural approaches include convolutional and recurrent networks \cite{roy2020} and transformer-based message embeddings combined with ensemble classifiers \cite{ghourabi2023}. On the Mendeley dataset, Mahmud et al. \cite{mahmud2024} combine a convolutional network with a bidirectional gated recurrent unit over word embeddings to separate legitimate from malicious messages, and Shim et al. \cite{shim2024} fine-tune BERT-family models, including BERT, DistilBERT and RoBERTa, to separate smishing from spam and study data augmentation with persuasion-based prompts. Further corpora include messages reported to Smishtank \cite{timko2024}, the Smishing-4C collection labelled with four smishing categories \cite{smishing4c2024}, MOZ-Smishing with Portuguese mobile-money fraud messages from Mozambique \cite{moz2025}, and the Super SMS Dataset, which was released together with a study of evasion techniques against spam filters \cite{salman2024}.

\emph{Character models and transformers.} Linear support vector machines over sparse n-gram features are strong text classifiers \cite{cortes1995,joachims1998}, and character n-grams are less sensitive to spelling variation than word features. Sigmoid (Platt) scaling turns the margins of such models into class probabilities \cite{niculescu2005}, which is needed for any rule that compares a confidence with a threshold; modern neural networks are often poorly calibrated as well \cite{guo2017}. Common baselines for short-text classification are random forests \cite{breiman2001}, gradient boosting \cite{chen2016}, convolutional networks \cite{kim2014} and long short-term memory networks \cite{hochreiter1997}. Pretrained transformer encoders \cite{devlin2019,sanh2019,liu2019} are the current standard for fine-tuned classification, and instruction-tuned large language models \cite{qwen2024} can classify messages from a few examples in the prompt.

\emph{Robustness to character-level perturbations.} DeepWordBug \cite{gao2018} and TextBugger \cite{li2019} generate small character and word edits that change the predictions of neural text classifiers. Pruthi et al. \cite{pruthi2019} defend against adversarial misspellings with a word recognition model placed before the classifier, and Boucher et al. \cite{boucher2022} describe perturbations with invisible characters, homoglyphs, reorderings and deletions. Salman et al. \cite{salman2024} evaluate evasion techniques against SMS spam filters. In this work, the perturbation types used for augmentation and those used only at test time are kept apart.

\emph{Cascades and early exit.} Cascades of classifiers of increasing cost go back at least to the detector of Viola and Jones \cite{viola2001}. In neural networks, early-exit methods stop the computation at an intermediate layer when the prediction is confident \cite{teerapittayanon2016,schwartz2020}. Schwartz et al. \cite{schwartz2020} use calibrated confidence to decide when to exit. Our cascade follows the same principle with two models of different kinds: a cheap character-level linear model and a transformer, with the exit decided by a calibrated probability and a threshold chosen on validation data.

\emph{Duplicates, leakage and the comparison of classifiers.} Leakage between training and test data, including duplicated records, is a common cause of over-optimistic results in applied machine learning \cite{kapoor2023}. Removing exact and near duplicates changes both training and evaluation of language models \cite{lee2022}. Near duplicates are commonly found by comparing sets of character shingles, with MinHash as an efficient approximation of their Jaccard similarity \cite{broder1997}. For the comparison of two classifiers on one test set, Dietterich \cite{dietterich1998} recommends the McNemar test \cite{mcnemar1947}; the bootstrap \cite{efron1993} gives confidence intervals for differences in metrics such as macro F1 that are not simple proportions; and the Holm procedure \cite{holm1979,aickin1996} controls the family-wise error rate when several comparisons are made.
